\uuid{5ZfL}
\titre{Somme de deux variables normales indépendantes}
\auteur{Maxime Nguyen}
\organisation{AMSCC}
\datecreate{2026-11-27}
\niveau{L2}
\module{Probabilité}
\chapitre{Probabilité continue}
\sousChapitre{Loi normale}
\difficulte{1}
\theme{loi normale, somme de variables indépendantes, variance}

\contenu{
\begin{question}
\qtype{onlyone}
On note $\mathcal{N}(\mu,\sigma)$ la loi normale d'espérance $\mu$ et d'écart-type $\sigma$. Soit $X_1 \sim \mathcal{N}(1, 3)$ et $X_2 \sim \mathcal{N}(2, 4)$, indépendantes. La loi de $X_1 + X_2$ est :
\begin{answers}
\good{$\mathcal{N}(3, 5)$}
\bad{$\mathcal{N}(3, 7)$}
\bad{$\mathcal{N}(3, 25)$}
\bad{$\mathcal{N}(3, \sqrt7)$}
\end{answers}
\begin{explanations}
Les espérances s'ajoutent toujours ; par indépendance,
ce sont les variances qui s'ajoutent : $9 + 16 = 25$, donc l'écart-type
vaut $\sqrt{25} = 5$.

B ajoute les écarts-types.

C donne la variance à la place de l'écart-type.

D : $\sqrt{3 + 4}$, les écarts-types n'ont pas été élevés au carré.

Contre-sens visé : additionner les écarts-types au lieu des variances.
\end{explanations}
\end{question}
}
